% Demonstration der Klasse tudcdthesis
%
% Für den Druck die Option output=print setzen. Dies aktiviert den doppelseitigen
% Satz mit Bindekorrektur und setzt alle farbigen Elemente schwarz.
\documentclass[
  output=digital,
  abstract={section,multiple,totoc},
  declaration={chapter,totoc},
]{tudcdthesis}
\usepackage[english,ngerman]{babel}
\usepackage[babel]{microtype}
\usepackage{amsmath}
\usepackage{csquotes}
\usepackage{blindtext}
\usepackage{hyperref}

% Organisationseinheiten
\faculty{Fakultät Kaffeemaschinenbau}
\institute{Institut für Typographie}
\chair{Professur für Angewandte Kompilierfehler}

% Angaben zur Arbeit
\thesis{diploma}
\graduation[Dipl.-Ing.]{Diplomingenieur}
\title{Über das ordentliche Setzen einer Abschlussarbeit}
\subtitle{im Corporate Design der TU~Dresden}
\author{Hans Wirdingenieur}
\matriculationnumber{1234567}
\dateofbirth{20.02.2002}
\placeofbirth{Dresden}
\course{Technische Textsatzkunde}
\discipline{Angewandte Gleitkommaverschiebung}
\supervisor{Dr.-Ing. Frieda Fließtext \and Dipl.-Ing. Moritz Marginpar}
\professor{Prof. Dr.-Ing. Berta Blocksatz}
\date{29. Februar 2026}

\begin{document}
\frontmatter
\maketitle

\begin{abstract}
  \blindtext
  \nextabstract[english]
  \blindtext
\end{abstract}

\begin{declarations}
  \confirmation
  \blocking[company=Heimlich GmbH]
\end{declarations}

\tableofcontents

\mainmatter
\chapter{Einleitung}
\blindtext

\section{Motivation}
\blindtext[2]

\subsection{Ein Unterabschnitt}
\blindtext

\paragraph{Ein Absatz}
\blindtext

\chapter{Grundlagen}
\section{Eine Gleichung}
Aus der Bewegungsgleichung~\eqref{eq:bewegung} folgt nach Abschnitt~\ref{sec:schluss}
das Ergebnis.
\begin{equation}
  m\ddot{x} + d\dot{x} + kx = F(t)
  \label{eq:bewegung}
\end{equation}
\blindtext

\section{Zusammenfassung}
\label{sec:schluss}
\blindtext

\end{document}
